\section{Vérifications numériques}
\label{sec:verifications_numeriques}

Cette section synthétise les vérifications numériques principales du
programme. Chaque ligne renvoie à la section qui en porte la
démonstration ou la description.

\subsection{Tableau récapitulatif}
\label{sec:verif_table}

\begin{longtable}{@{}p{6.5cm}p{4.5cm}p{3cm}@{}}
\toprule
\textbf{Test} & \textbf{Précision / score} & \textbf{Référence}\\
\midrule
\endhead
Identité Fredholm $\det(I - \DRop)^{-1} = \Rres(s)$, $\Reop(s) > 1$
 & $\sim 2 \cdot 10^{-12}$ & \S\ref{sec:fredholm}\\
Formule de trace $-d\log \Rres/ds = \sum_p (\log p)[1/(p^s-1) + \Ap p^{-s}]$
 & $\sim 10^{-26}$ & thm.~\ref{thm:trace_formula}\\
Différence $\NPT - N_{\mathrm{RvM}} = 1/8$ constante sur 4 décades
 & $\sim 10^{-15}$ & \S\ref{sec:densite_semi_classique}\\
Symétrie spectrale $\{\gamma_n, -\gamma_n\}$ (BC AP, $\theta = \pi$)
 & $5 \cdot 10^{-12}$ & prop.~\ref{prop:antiperiodicite}\\
Phase de Berry par coin $-\pi/4$, densité $-1/8$
 & $1{,}68 \cdot 10^{-14}$ & prop.~\ref{prop:c1_sur_N}\\
$\epsstar^{\mathrm{geom}}$ vs optimum empirique $W = 24$
 & $\sim 5\,\%$ global, $\sim 1\,\%$ au régime $M = 1155, W = 24$
 & lem.~\ref{lem:epsstar_geom}\\
Pôles de $|\Treg(1/2 + it)|$ (formule brute) sur $t \in [10, 70]$
 & $17/17$ à $|\Delta| < 1{,}5$
 & \S\ref{sec:trace_reg}\\
Pôles de $|\Treg(1/2 + it)|$ (formule brute) sur $t \in [10, 200]$
 & $74/79$ à $|\Delta| < 1{,}5$ (sat.~$P_{\max} = 3\,000$)
 & \S\ref{sec:trace_reg}\\
Pôles de $|\Treg(1/2 + it)|$ avec $R_1(\eps = 0{,}2)$ sur $t \in [10, 200]$
 & $77/79$ à $|\Delta| < 1{,}5$, \textbf{$75/79$ à $|\Delta| < 0{,}5$}, $47/79$ à $|\Delta| < 0{,}1$
 & \S\ref{sec:trace_reg}\\
Pôles de $|\Treg(1/2 + it)|$ avec $R_1(\eps = 0{,}2)$ sur $t \in [10, 500]$
 & $267/269$ ($99{,}3\,\%$), \textbf{$242/269$ ($90{,}0\,\%$) à $|\Delta| < 0{,}5$}, $139/269$ ($51{,}7\,\%$) à $|\Delta| < 0{,}1$
 & \S\ref{sec:verif_smearing}\\
\bottomrule
\caption{Vérifications numériques principales du programme HP-PT.}
\label{tab:verifs_numeriques}
\end{longtable}

\subsection{Forme fonctionnelle exacte de $\NPT$}
\label{sec:verif_NPT}

Vérification de \eqref{eq:NPT} sur
$\gamma \in \{14, 21, 30, 50, 100, 200, 500, 1000\}$~:

\begin{table}[ht]
\centering
\begin{tabular}{@{}rrrr@{}}
\toprule
$\gamma$ & $\NPT$ & $N_{\mathrm{RvM}}$ & différence\\
\midrule
$14{,}13$ & $0{,}5737$ & $0{,}4487$ & $0{,}1250$\\
$21{,}02$ & $1{,}6945$ & $1{,}5695$ & $0{,}1250$\\
$50{,}00$ & $9{,}5478$ & $9{,}4228$ & $0{,}1250$\\
$100$ & $29{,}1273$ & $29{,}0023$ & $0{,}1250$\\
$500$ & $269{,}7117$ & $269{,}5867$ & $0{,}1250$\\
$1\,000$ & $648{,}7412$ & $648{,}6162$ & $0{,}1250$\\
\bottomrule
\end{tabular}
\caption{Différence $\NPT - N_{\mathrm{RvM}}$ constante à $1/8$ sur quatre
décades, à précision machine ($\le 10^{-15}$).}
\label{tab:verif_residu}
\end{table}

\subsection{Capture distributionnelle des zéros}
\label{sec:verif_capture}

Test sur $t \in [10, 70]$ avec régularisation Abel--Hadamard $p^{-\eps}$.
Pour $\eps \in \{0{,}6,\,0{,}4,\,0{,}2\}$, le module $|F_\eps(t)|$ croît
systématiquement de $2{,}7\times$ à $3{,}3\times$ aux zéros de $\zeta$
quand $\eps$ passe de $0{,}6$ à $0{,}2$ (signature d'un pôle), contre
$1{,}4\times$ à $2{,}3\times$ aux points neutres (régularité). Sur
$t \in [10, 70]$, $dt = 0{,}1$, on détecte 41-52 pics locaux selon
$P_{\max} \in [3\,000,\, 100\,000]$ ; les 17 zéros de référence sont
tous à moins de $1{,}5$ unité d'un pic ($|\Delta| < 1{,}5$) dès
$P_{\max} = 3\,000$.

\paragraph{Saturation $P_{\max}$.}
La capture sature au-delà de $P_{\max} = 3\,000$~: la médiane $|\Delta|$
décroît de $0{,}49$ ($P_{\max} = 3\,000$) à $0{,}36$ ($P_{\max} = 30\,000$)
puis stagne. La même statistique étendue à $t \in [10, 200]$ donne
$74/79$ zéros (94\,\%) à $|\Delta| < 1{,}5$, $38$ à $|\Delta| < 0{,}5$,
$7$ à $|\Delta| < 0{,}1$. Les 5 zéros manqués appartiennent à des paires
$(\gamma_n, \gamma_{n+1})$ avec écart inférieur à la largeur intrinsèque
d'un pic non régularisé ($\sim 2\pi/\log\gamma$), fusionnées en un pic
unique. Cas extrême de précision~: $\gamma = 146{,}001$ capturé à
$|\Delta| = 0{,}001$. L'amélioration sub-$0{,}1$ systématique requiert
d'appliquer la régularisation explicite $R_1$ ou $R_2$, et non
l'augmentation de $P_{\max}$.

\subsection{Régularisation par smearing gaussien et range étendu}
\label{sec:verif_smearing}

Avec $\phi_\eps(t) = (\eps \sqrt{2\pi})^{-1} \exp(-t^{2} / 2\eps^{2})$ et
$\eps \in \{2,\,1,\,0{,}5\}$, le contraste pôle/neutre est plus marqué~:
$8\times$ à $10\times$ aux zéros contre $2\times$ à $3\times$ aux points
neutres. La régularisation gaussienne et la régularisation Abel sont
équivalentes au sens distributionnel (différence de noyaux régularisants),
mais la première est plus sensible numériquement.

\paragraph{Range étendu $t \in [10, 500]$ (269 zéros).}
$R_1$ avec $\eps = 0{,}2$ capture $267/269$ zéros à $|\Delta| < 1{,}5$,
$242/269$ à $|\Delta| < 0{,}5$ et $139/269$ à $|\Delta| < 0{,}1$. La
décomposition par bande de 100 unités est~:

\begin{table}[ht]
\centering
\begin{tabular}{@{}lrr@{}}
\toprule
\textbf{Bande} & \textbf{Zéros} & \textbf{Capture $|\Delta| < 0{,}5$}\\
\midrule
$[10,\,100]$ & 38 & $100{,}0\,\%$\\
$[100,\,200]$ & 41 & $92{,}0\,\%$\\
$[200,\,300]$ & 50 & $89{,}8\,\%$\\
$[300,\,400]$ & 67 & $87{,}5\,\%$\\
$[400,\,500]$ & 73 & $86{,}6\,\%$\\
\bottomrule
\end{tabular}
\caption{Capture par sous-bande de $100$ unités, $R_1$ smearing gaussien
$\eps = 0{,}2$. Dégradation monotone mais bornée : la performance ne
s'effondre pas avec $T$.}
\label{tab:bandes}
\end{table}

Le plancher à $86{,}6\,\%$ est cohérent avec la limite structurelle de
$55$ paires fusionnées (gap $< 2\pi/\log\gamma$) sur $269$ zéros, soit
$20{,}4\,\%$ des zéros appartenant à une paire au-dessous de la
résolution intrinsèque du modèle.

\paragraph{Stratégie $\eps$-adaptative testée.}
Une stratégie $\eps$-adaptative ($\eps \propto 1/\log\gamma$) a été
testée sous deux formes~: convolution sur $|F|$ brut (effondrement des
pics, falsifiée) et convolution sur $F$ complexe (compromis marginal~:
médiane $|\Delta|$ réduite à $0{,}080$ mais couverture réduite à
$76\,\%$ à $|\Delta| < 1{,}5$). La régularisation à poids spectral
fixe ($R_1$, $\eps = 0{,}2$) reste l'\textbf{optimum global}.

